\documentclass[10pt]{article}

% ====================================================
% PACKAGES
% ====================================================
\usepackage[margin=0.5in, headheight=13.6pt, columnsep=20pt]{geometry}
\usepackage[T1]{fontenc}
\usepackage{lmodern}
\usepackage{microtype}
\usepackage{setspace}
\usepackage{tcolorbox}
\usepackage{amsmath, amssymb, mathtools}
\usepackage{siunitx}

\usepackage{graphicx}
\usepackage{float}
\usepackage{subcaption}

\usepackage{enumitem}
\usepackage{booktabs}
\usepackage{longtable}
\usepackage{array}
\usepackage{multirow}

\usepackage{xcolor}
\usepackage{hyperref}

\usepackage{fancyhdr}
\usepackage{lastpage}

\usepackage{listings}
\usepackage{indentfirst}
\usepackage[justification=centering]{caption}

\usepackage{multicol}
\usepackage{titlesec}

% Python source style for \CodeListing
\lstdefinestyle{pystyle}{
    language=Python,
    basicstyle=\ttfamily\footnotesize,
    keywordstyle=\color{blue!60!black}\bfseries,
    stringstyle=\color{orange!70!black},
    commentstyle=\color{gray!60!black}\itshape,
    numbers=left,
    numberstyle=\tiny\color{gray!70},
    numbersep=8pt,
    breaklines=true,
    showstringspaces=false,
    frame=single,
    rulecolor=\color{gray!40},
    tabsize=4,
    columns=fullflexible,
    captionpos=b,
    literate={°}{{\textdegree}}1,
}
\lstset{style=pystyle}

\pagestyle{fancy}
\fancyhf{}
\lhead{\Assignment\ |\ \course}
\rhead{\Name\ |\ \today}
\cfoot{Page \thepage\ of~\pageref{LastPage}}

\hypersetup{
    colorlinks,
    linkcolor={red!50!black},
    citecolor={blue!50!black},
    urlcolor={blue!80!black}
}

\newcommand{\Name}{Arael Anaya}
\newcommand{\Assignment}{Equation Sheet}
\newcommand{\course}{Advanced Dynamics}

% Boxed key-formula environment (tight padding for cheat-sheet density)
\newtcolorbox{keybox}{
    colback=blue!4!white,
    colframe=blue!45!black,
    boxrule=0.5pt,
    arc=1pt,
    left=4pt, right=4pt, top=1pt, bottom=1pt,
    before skip=2pt, after skip=2pt
}

% Compact section/subsection formatting
\titleformat{\section}{\normalfont\bfseries\large\color{blue!40!black}}{}{0em}{}[{\color{blue!40!black}\titlerule[0.8pt]}]
\titlespacing*{\section}{0pt}{4pt}{1pt}
\titleformat{\subsection}{\normalfont\bfseries\color{blue!25!black}}{}{0em}{}
\titlespacing*{\subsection}{0pt}{2pt}{1pt}

\setlength{\columnseprule}{0.4pt}
\def\columnseprulecolor{\color{gray!40}}
\setlength{\multicolsep}{2pt}
\setlength{\parskip}{0pt}

% Notation shorthand
\newcommand{\rate}[2]{{}^{#1}\dfrac{d}{dt}#2}

\begin{document}
\thispagestyle{fancy}
\begin{center}
    {\Large\bfseries Advanced Dynamics --- Equation Sheet}\\[1pt]
    {\small Kinematics \(\cdot\) Newton's Laws \(\cdot\) Energy \& Momentum \(\cdot\) Generalized Coordinates \& Constraints \(\cdot\) D'Alembert's Principle \(\cdot\) Lagrangian \& Hamiltonian Mechanics \(\cdot\) Nonlinear Dynamics \& Bifurcations \(\cdot\) Nondimensionalization}
\end{center}
\vspace{-8pt}

\begin{multicols}{2}
\footnotesize

% ====================================================
% NOTATION
% ====================================================
\section*{Notation}
\begin{itemize}[nosep, leftmargin=1.1em]
    \item $N, A, B,\dots$ --- reference frames (e.g.\ $N$ = fixed/inertial)
    \item $\hat n_i, \hat a_i, \hat b_i$ --- right-handed unit-vector basis fixed in $N,A,B$
    \item $\vec r_{P/O}$ --- position of point $P$ relative to point $O$; $\vec r\equiv{}^{N}\vec r$ unless noted otherwise
    \item ${}^{N}\vec v_P,\ {}^{N}\vec a_P$ --- velocity/acceleration of $P$ \emph{as seen from} $N$
    \item $\vec\omega^{A/N}$ --- angular velocity of frame $A$ relative to $N$
    \item $\vec\alpha^{A/N} = {}^{N}\dfrac{d}{dt}\vec\omega^{A/N}$ --- angular acceleration of $A$ relative to $N$
    \item $\rate{N}{(\cdot)}$ --- time derivative taken \emph{by an observer fixed in} $N$
    \item $[R]^{B/A}$ --- rotation (direction cosine) matrix: $\{\hat b\} = [R]^{B/A}\{\hat a\}$
    \item $q_1,\dots,q_N$ --- generalized coordinates; $\vec\gamma_{ij}=\partial\vec r_i/\partial q_j$ --- partial velocities
    \item $\vec P = m\vec v$ --- linear momentum; $\vec H_O=\vec r\times\vec P$ --- angular momentum about $O$
    \item $L=T-V$ --- Lagrangian; $p_j=\partial L/\partial\dot q_j$ --- conjugate (generalized) momentum
    \item $H=\sum_jp_j\dot q_j-L$ --- Hamiltonian; $\lambda_k$ --- Lagrange multiplier for constraint $\phi_k$
\end{itemize}

% ====================================================
% ROTATIONS & ANGULAR VELOCITY
% ====================================================
\section*{Rotations \& Angular Velocity}

\subsection*{Rotation matrices}
Orthogonal, so $[R]^{-1} = [R]^{T}$; composing rotations through an intermediate frame $A$:
\[
    [R]^{B/N} = [R]^{B/A}[R]^{A/N}
\]

\subsection*{Angular velocity addition (chain rule)}
Valid for any chain of intermediate frames:
\begin{keybox}
\[
    \vec\omega^{B/N} = \vec\omega^{B/A} + \vec\omega^{A/N},
    \qquad
    \vec\omega^{C/N} = \vec\omega^{C/B}+\vec\omega^{B/A}+\vec\omega^{A/N}
\]
\end{keybox}
Note the reversal/cancellation identities:
\[
    \vec\omega^{A/N} = -\vec\omega^{N/A}, \qquad \vec\omega^{A/A} = \vec 0
\]

\subsection*{Derivative of a rotating unit vector}
If $\hat b_i$ is fixed in $B$ (constant length/orientation relative to $B$):
\[
    \rate{N}{\hat b_i} = \vec\omega^{B/N}\times\hat b_i
\]

\subsection*{2D rotation example ($N,P$ frames offset by $\alpha$)}
\[
    \hat n_1 = \hat p_1\cos\alpha+\hat p_2\sin\alpha, \qquad \hat n_2=-\hat p_1\sin\alpha+\hat p_2\cos\alpha
\]
\[
    \begin{Bmatrix}\hat n_1\\\hat n_2\end{Bmatrix} = [R]^{N/P}\begin{Bmatrix}\hat p_1\\\hat p_2\end{Bmatrix},
    \qquad
    [R]^{N/P} = \begin{bmatrix}\cos\alpha&\sin\alpha\\-\sin\alpha&\cos\alpha\end{bmatrix}
\]
(written ${}^{N}R_P$ in lecture)
\[
    [R]^{P/N} = \big([R]^{N/P}\big)^{-1} = \big([R]^{N/P}\big)^{T} = \begin{bmatrix}\cos\alpha&-\sin\alpha\\\sin\alpha&\cos\alpha\end{bmatrix}
\]

% ====================================================
% VECTOR DIFFERENTIATION FORMULA (VDF)
% ====================================================
\section*{Vector Differentiation Formula (VDF)}
Also called the \emph{transport theorem}: for \emph{any} vector $\vec q$ (not necessarily fixed in either frame), relating the rate of change seen in $N$ to the rate of change seen in a frame $A$ rotating at $\vec\omega^{A/N}$:
\begin{keybox}
\[
    \rate{N}{\vec q} = \rate{A}{\vec q} + \vec\omega^{A/N}\times\vec q
\]
\end{keybox}
\textbf{Corollary:} if $\vec q$ is fixed in $A$ (i.e.\ $\rate{A}{\vec q}=0$), this reduces to $\rate{N}{\vec q}=\vec\omega^{A/N}\times\vec q$ (previous box). This is the tool used to differentiate a vector expressed in a rotating basis one component at a time.

% ====================================================
% POSITION VECTOR ADDITION
% ====================================================
\section*{Position Vector Addition (Setup for Moving Frames)}
For two points $P$ and $O'$ both located relative to a common origin $O$ (equivalently $\vec r=\vec R+\vec r\,'$, with $\vec R=\vec r_{O'/O}$, $\vec r\,'=\vec r_{P/O'}$), position vectors add/subtract exactly like any vectors --- the basic relation used to move between frames whose origins differ:
\begin{keybox}
\[
    \vec r_{P/O} = \vec r_{P/O'} + \vec r_{O'/O}
    \qquad\Longleftrightarrow\qquad
    \vec r_{P/O'} = \vec r_{P/O} - \vec r_{O'/O}
\]
\end{keybox}
Chained over several intermediate points: $\vec r_{P/O} = \vec r_{P/B} + \vec r_{B/A} + \vec r_{A/O}$. Applying the VDF to this relation term-by-term is exactly how the velocity/acceleration relations below are derived.

% ====================================================
% TRANSLATING & ROTATING FRAMES --- VELOCITY
% ====================================================
\section*{Velocity: Translating \& Rotating Frames}
Let $O'$ be the (moving) origin of frame $A$, which translates \emph{and} rotates ($\vec\omega^{A/N}$) relative to $N$, and let $P$ be an arbitrary point:
\begin{keybox}
\[
    {}^{N}\vec v_P = {}^{N}\vec v_{O'} + {}^{A}\vec v_P + \vec\omega^{A/N}\times\vec r_{P/O'}
\]
\end{keybox}
\begin{itemize}[nosep, leftmargin=1.1em]
    \item ${}^{N}\vec v_{O'}$ --- translational velocity of the moving origin
    \item ${}^{A}\vec v_P = \rate{A}{\vec r_{P/O'}}$ --- velocity of $P$ \emph{relative to} $A$ (as seen by an observer riding on $A$)
    \item $\vec\omega^{A/N}\times\vec r_{P/O'}$ --- velocity contribution from $A$'s rotation
\end{itemize}
\textbf{Special cases}
\begin{itemize}[nosep, leftmargin=1.1em]
    \item Pure translation ($\vec\omega^{A/N}=0$): ${}^{N}\vec v_P = {}^{N}\vec v_{O'} + {}^{A}\vec v_P$
    \item $P$ fixed in $A$ (${}^{A}\vec v_P = 0$, rigid-body point): ${}^{N}\vec v_P = {}^{N}\vec v_{O'} + \vec\omega^{A/N}\times\vec r_{P/O'}$
    \item Common origin ($O'=O$, ${}^{N}\vec v_{O'}=0$): reduces to the transport theorem applied to $\vec r_{P/O}$
\end{itemize}

% ====================================================
% TRANSLATING & ROTATING FRAMES --- ACCELERATION
% ====================================================
\section*{Acceleration: Translating \& Rotating Frames}
Differentiating the velocity relation again (the ``five-term'' / Coriolis theorem):
\begin{keybox}
\begin{align*}
    {}^{N}\vec a_P ={}& {}^{N}\vec a_{O'} + {}^{A}\vec a_P
    + \vec\alpha^{A/N}\times\vec r_{P/O'} \\
    &+ \vec\omega^{A/N}\times\big(\vec\omega^{A/N}\times\vec r_{P/O'}\big)
    + 2\,\vec\omega^{A/N}\times {}^{A}\vec v_P
\end{align*}
\end{keybox}
\begin{itemize}[nosep, leftmargin=1.1em]
    \item ${}^{N}\vec a_{O'}$ --- translational acceleration of the moving origin
    \item ${}^{A}\vec a_P = \rate{A}{{}^{A}\vec v_P}$ --- acceleration of $P$ relative to $A$
    \item $\vec\alpha^{A/N}\times\vec r_{P/O'}$ --- Euler (angular acceleration) term
    \item $\vec\omega^{A/N}\times(\vec\omega^{A/N}\times\vec r_{P/O'})$ --- centripetal term
    \item $2\,\vec\omega^{A/N}\times {}^{A}\vec v_P$ --- Coriolis term
\end{itemize}
\textbf{Special cases}
\begin{itemize}[nosep, leftmargin=1.1em]
    \item Pure translation ($\vec\omega^{A/N}=0$): ${}^{N}\vec a_P = {}^{N}\vec a_{O'} + {}^{A}\vec a_P$
    \item $P$ fixed in $A$ (${}^{A}\vec v_P={}^{A}\vec a_P=0$, rigid-body point):
    \[
        {}^{N}\vec a_P = {}^{N}\vec a_{O'} + \vec\alpha^{A/N}\times\vec r_{P/O'} + \vec\omega^{A/N}\times(\vec\omega^{A/N}\times\vec r_{P/O'})
    \]
    \item $P$ moving in $A$ but $O'$ fixed in $N$ (common origin, no translation): drop the ${}^{N}\vec a_{O'}$ term --- this is the form used for rotating-basis problems (e.g.\ spherical pendulum, rotating tube)
\end{itemize}

% ====================================================
% NEWTON'S LAWS
% ====================================================
\section*{Newton's Laws of Motion}
\begin{keybox}
\[
    \textbf{1st:}\ \textstyle\sum\vec F=0 \Rightarrow {}^{N}\ddot{\vec r}=\text{const}
    \qquad
    \textbf{2nd:}\ \textstyle\sum\vec F = m\,{}^{N}\ddot{\vec r}
\]
\end{keybox}
\[
    \textbf{3rd:}\ \ \vec F_{12} = -\vec F_{21}\quad\text{(action--reaction pairs)}
\]
Newton's 2nd law \textbf{requires} the acceleration relative to an inertial (Newtonian) frame $N$ --- never a rotating/accelerating frame directly.

% ====================================================
% WORK & ENERGY
% ====================================================
\section*{Work \& Energy}
\subsection*{Kinetic \& potential energy}
\[
    T=\tfrac12 m|\dot{\vec r}|^{2}=\tfrac12 m\,\dot{\vec r}\cdot\dot{\vec r},
    \qquad
    V=V(\vec r)
\]
\[
    \vec F=-\nabla V(\vec r)\quad\text{(conservative/``potential'' force)}
\]
\textbf{Common potentials}
\begin{itemize}[nosep, leftmargin=1.1em]
    \item Uniform gravity: $V=mgh$, $\vec F=-mg\,\hat n_3$
    \item Newtonian gravity: $V=-\dfrac{GmM}{r}$, $\vec F=\dfrac{GmM}{r^{2}}\hat b_1$ ($\hat b_1$: unit vector toward the attracting mass)
    \item Linear spring (stretch $x$): $V=\tfrac12 kx^{2}$, $\vec F=-kx\,\hat n_1$
\end{itemize}
\subsection*{Work--energy theorem}
\[
    W=\int_{\vec r_1}^{\vec r_2}\vec F\cdot d\vec r, \qquad T_2-T_1=W
\]
\begin{keybox}
All forces conservative ($\vec F=-\nabla V$):
\[
    T_1+V_1=T_2+V_2 \quad(\text{energy conserved})
\]
\end{keybox}
Only holds when \emph{every} force doing work is conservative --- friction or other non-conservative applied forces break it.

% ====================================================
% MOMENTUM
% ====================================================
\section*{Momentum}
\subsection*{Linear momentum}
\[
    \vec P=m\dot{\vec r}, \qquad \dot{\vec P}=m\ddot{\vec r}=\textstyle\sum\vec F
\]
If $\sum\vec F=\vec 0$, $\vec P$ is conserved.
\subsection*{Angular momentum (about inertial fixed point $O$)}
\begin{keybox}
\[
    \vec H_O=\vec r\times\vec P, \qquad \dot{\vec H}_O=\vec r\times\textstyle\sum\vec F=\textstyle\sum\vec M_O
\]
\end{keybox}
``Rotational Newton's 2nd law.'' If $\sum\vec M_O=\vec 0$, $\vec H_O$ is conserved.

% ====================================================
% GENERALIZED COORDINATES
% ====================================================
\section*{Generalized Coordinates}
\begin{itemize}[nosep, leftmargin=1.1em]
    \item \textbf{Generalized coordinates} $q_1,\dots,q_N$: minimal (or redundant) geometric variables (angles/distances) describing the configuration of a system
    \item \textbf{Configuration space:} the set of all configurations the system can take on, parametrized by the $q_j$
    \item A system may be described by a \textbf{redundant set} (more coordinates than DOF, related by constraint equations) or a \textbf{minimal set} (one coordinate per DOF, no constraints left)
\end{itemize}
\begin{keybox}
\[
    n_{\text{DOF}} = \#(\text{gen.\ coords}) - \#(\text{indep.\ constraints})
\]
\end{keybox}
Always reduce to the minimal set before deriving equations of motion.

% ====================================================
% HOLONOMIC CONSTRAINTS
% ====================================================
\section*{Holonomic Constraints}
\textbf{Position-based}: with $N$ generalized coordinates and $M$ constraints ($j=1,\dots,M$):
\begin{keybox}
\[
    \phi_j(q_1,\dots,q_N)=\phi_j(\vec q)=0, \quad\text{or, if time-dependent,}\quad \phi_j(\vec q,t)=0
\]
\end{keybox}
\begin{itemize}[nosep, leftmargin=1.1em]
    \item \textbf{Scleronomic:} no explicit $t$ --- time-independent constraint
    \item \textbf{Rheonomic:} explicit $t$ dependence --- e.g.\ a prescribed moving support
    \item Each independent holonomic constraint removes exactly one DOF (box above)
    \item \textbf{Always rewritable as non-holonomic:} since $\phi_j\equiv0\Rightarrow\dot\phi_j\equiv0$,
    \[
        \dot\phi_j(\vec q,t)=\sum_{i=1}^{N}\frac{\partial\phi_j}{\partial q_i}\dot q_i+\frac{\partial\phi_j}{\partial t}=0
        \qquad\big(a_{ji}\equiv\partial\phi_j/\partial q_i\big)
    \]
\end{itemize}

% ====================================================
% NON-HOLONOMIC CONSTRAINTS
% ====================================================
\section*{Non-Holonomic Constraints}
\textbf{Velocity-based}: $\phi_j(\vec q,\dot{\vec q},t)=0$, usually taken \emph{linear} in the velocities (Pfaffian form):
\begin{keybox}
\[
    \phi_j=\sum_{i=1}^{N}a_{ji}(\vec q,t)\,\dot q_i+a_{jt}(\vec q,t)=0
\]
\end{keybox}
(no explicit $t$: $\ \phi_j=\sum_i a_{ji}(\vec q)\dot q_i=0$)
\begin{itemize}[nosep, leftmargin=1.1em]
    \item Every holonomic constraint can be written this way (box above); the \textbf{converse is false} --- a Pfaffian constraint is genuinely non-holonomic unless it is exact/integrable (or made exact with an integrating factor; check via the Frobenius condition)
    \item Non-holonomic constraints reduce the number of independent \emph{velocities} but \textbf{not} the dimension of the configuration space
    \item \textbf{Example --- knife-edge/skate/wheel constraint:} contact-point velocity must be perpendicular to the blade, $\vec v_A\cdot\hat n=0$ with $\hat n=(-\sin\theta,\cos\theta)$:
    \[
        \phi=-\sin\theta\,\dot x+\cos\theta\,\dot y=0
    \]
    With $(q_1,q_2,q_3)=(x,y,\theta)$: $a_1=-\sin q_3,\ a_2=\cos q_3,\ a_3=0$. Cannot be integrated to $\phi(\vec q)=0$ --- genuinely non-holonomic.
\end{itemize}

% ====================================================
% D'ALEMBERT'S PRINCIPLE
% ====================================================
\section*{D'Alembert's Principle}
\subsection*{Partial velocities}
For particle $i$ with position $\vec r_i(\vec q,t)$, the partial velocity w.r.t.\ generalized speed $\dot q_j$:
\[
    \vec\gamma_{ij}=\frac{\partial\vec r_i}{\partial q_j}=\frac{\partial\dot{\vec r}_i}{\partial\dot q_j}
\]
(equal because $\dot{\vec r}_i$ is linear in the $\dot q_j$'s --- fastest computed by differentiating the already-known $\dot{\vec r}_i$ directly w.r.t.\ $\dot q_j$.)

\subsection*{Equations of motion}
For a minimal set $q_1,\dots,q_M$ and particles $i=1,\dots,n$:
\begin{keybox}
\[
    \sum_{i=1}^{n}\big(\vec F_{Ai}-m_i\ddot{\vec r}_i\big)\cdot\vec\gamma_{ij}=0, \qquad j=1,\dots,M
\]
\end{keybox}
\begin{itemize}[nosep, leftmargin=1.1em]
    \item $\vec F_{Ai}$ --- \emph{applied} (impressed) forces on particle $i$ only (gravity, springs, motors, \dots)
    \item \textbf{Ideal constraint forces are excluded} (rod tension, normal forces, frictionless-contact forces) --- they do no virtual work consistent with the constraints and drop out automatically
    \item One scalar equation per generalized coordinate --- exactly $M$ equations for $M$ DOF
    \item Equivalently, with generalized force $Q_j=\sum_i\vec F_{Ai}\cdot\vec\gamma_{ij}$: $\ \sum_i m_i\ddot{\vec r}_i\cdot\vec\gamma_{ij}=Q_j$
\end{itemize}

\subsection*{Recipe}
\begin{enumerate}[nosep, leftmargin=1.3em]
    \item Choose a minimal set $q_1,\dots,q_M$ consistent with all holonomic constraints
    \item Write $\vec r_i(\vec q,t)$ for every particle; differentiate for $\dot{\vec r}_i,\ddot{\vec r}_i$
    \item Compute $\vec\gamma_{ij}=\partial\dot{\vec r}_i/\partial\dot q_j$ for every particle/coordinate pair
    \item Identify applied forces $\vec F_{Ai}$ on each particle (\emph{exclude} constraint forces)
    \item For each $j$: $\sum_i(\vec F_{Ai}-m_i\ddot{\vec r}_i)\cdot\vec\gamma_{ij}=0$ $\Rightarrow$ $M$ equations of motion
    \item For linear (spring/mass) systems, these assemble into $[M]\ddot{\vec q}+[K]\vec q=\vec F$
\end{enumerate}

% ====================================================
% LAGRANGIAN MECHANICS
% ====================================================
\section*{Lagrangian Mechanics}
\subsection*{Euler--Lagrange equations}
For a minimal set $q_1,\dots,q_M$; equivalent to D'Alembert's principle above (substitute $\vec\gamma_{ij}=\partial\dot{\vec r}_i/\partial\dot q_j$ into $L=T-V$ and simplify):
\begin{keybox}
\[
    \frac{d}{dt}\left(\frac{\partial L}{\partial\dot q_j}\right) - \frac{\partial L}{\partial q_j} = Q_j^{nc}, \qquad j=1,\dots,M
\]
\end{keybox}
\begin{itemize}[nosep, leftmargin=1.1em]
    \item $Q_j^{nc}$ --- generalized force from any applied force \emph{not} already folded into $V$ (a prescribed drive $F(t)$, friction, \dots); $=0$ if every applied force is conservative
    \item One 2nd-order ODE per DOF; ideal constraint forces are still automatically excluded, since $L(\vec q,\dot{\vec q},t)$ is built from the minimal coordinates
    \item Cross terms quadratic in velocity (Coriolis-type, from $\partial T/\partial q_j$) fall out of $\tfrac{d}{dt}(\partial L/\partial\dot q_j)-\partial L/\partial q_j$ automatically --- no need to track them separately as in Newtonian/D'Alembert bookkeeping
\end{itemize}

\subsection*{Matrix form \& small oscillations}
Multi-DOF equations of motion assemble into
\[
    M(\vec q)\ddot{\vec q} + \vec C(\vec q,\dot{\vec q}) + \vec K(\vec q,t) = \vec 0
\]
($M$: symmetric positive-definite mass/inertia matrix; $\vec C$: Coriolis/centripetal-type $\dot q^2$ terms; $\vec K$: stiffness/gravity/forcing terms). Linearizing about an equilibrium ($\sin(\cdot)\to(\cdot)$, $\cos(\cdot)\to1$, drop $\dot q^2$ terms):
\begin{keybox}
\[
    M\ddot{\vec q} + K\vec q = \vec 0
\]
\end{keybox}
Assume $\vec q=\vec Qe^{i\omega t}$: non-trivial solutions require the \textbf{characteristic equation}
\[
    \det\!\big(K-\omega^{2}M\big)=0 \quad\Rightarrow\quad \text{natural frequencies }\omega_i
\]
Mode shape $\vec Q_i$ (up to overall scale) solves $(K-\omega_i^2M)\vec Q_i=\vec 0$; lowest $\omega_i$ = coordinates in phase, higher $\omega_i$ = out of phase.

\subsection*{Ignorable (cyclic) coordinates \& Routhian reduction}
If $L$ has no explicit $q_k$-dependence ($\partial L/\partial q_k=0$; only $\dot q_k$ appears), Lagrange's equation for it collapses to a conservation law:
\begin{keybox}
\[
    \frac{\partial L}{\partial\dot q_k} = B_k = \text{const} \qquad\text{(conserved generalized momentum)}
\]
\end{keybox}
Invert for $\dot q_k=f(\vec q_{\text{rest}},B_k)$ and use it to eliminate $\dot q_k$ from the remaining dynamics via the \textbf{Routhian}
\[
    R(\vec q_{\text{rest}},\dot{\vec q}_{\text{rest}};B_k) \equiv L - \dot q_k B_k\Big|_{\dot q_k\to f(\vec q_{\text{rest}},B_k)}
\]
\[
    \frac{d}{dt}\frac{\partial R}{\partial\dot q_j}-\frac{\partial R}{\partial q_j}=0 \quad (j\neq k)
\]
This reduces the problem by one coordinate; recover $q_k(t)$ afterward by integrating $\dot q_k=f(\vec q_{\text{rest}}(t),B_k)$.

\subsection*{Constraint forces via Lagrange multipliers}
For $M$ coordinates with $m$ Pfaffian constraints $\phi_k=\sum_ia_{ki}\dot q_i+a_{kt}=0$ ($k=1,\dots,m$; a holonomic constraint left un-eliminated works the same way, with $a_{ki}=\partial\phi_k/\partial q_i$):
\begin{keybox}
\[
    \frac{d}{dt}\frac{\partial L}{\partial\dot q_j}-\frac{\partial L}{\partial q_j} = \sum_{k=1}^{m}\lambda_k\,a_{kj}, \qquad j=1,\dots,M
\]
\end{keybox}
\begin{itemize}[nosep, leftmargin=1.1em]
    \item $M$ equations of motion $+$ $m$ constraint equations $\phi_k=0$ close the system for $\vec q(t)$ \textbf{and} $\lambda_k(t)$
    \item $\sum_k\lambda_k a_{kj}$ is exactly the constraint force conjugate to $q_j$ --- this is how the forces D'Alembert/plain Lagrange discard get recovered (e.g.\ rod tension, knife-edge normal force)
    \item \textbf{Recipe:} write $L=T-V$; list $a_{ki}$ from each $\phi_k$; write the $M$ boxed equations plus the $m$ constraints $\phi_k=0$; solve algebraically for $\lambda_k$ (introducing a body-frame speed like $u\equiv\dot{\vec r}\cdot\hat b_1$ often decouples the algebra)
\end{itemize}

% ====================================================
% HAMILTONIAN MECHANICS
% ====================================================
\section*{Hamiltonian Mechanics}
\subsection*{Generalized momentum \& the Hamiltonian}
\[
    p_j \equiv \frac{\partial L}{\partial\dot q_j}
\]
\begin{keybox}
\[
    H(\vec q,\vec p,t) \equiv \sum_j p_j\dot q_j - L \qquad\text{(Legendre transform of }L\text{)}
\]
\end{keybox}
Invert $p_j(\vec q,\dot{\vec q},t)$ for $\dot q_j(\vec q,\vec p,t)$ and substitute so $H$ is expressed purely in $(\vec q,\vec p,t)$.

\subsection*{When does $H=T+V$?}
Split $T$ by degree in the $\dot q_j$: $T=T_2+T_1+T_0$ ($T_2$ homogeneous quadratic, $T_1$ linear, $T_0$ velocity-independent --- all can depend on $\vec q,t$). In general
\begin{keybox}
\[
    H = T_2-T_0+V \qquad\text{(the \emph{Jacobi integral})}
\]
\end{keybox}
\begin{itemize}[nosep, leftmargin=1.1em]
    \item Reduces to $H=T+V$ \textbf{only} when $T_1=T_0=0$: $\vec r_i(\vec q)$ has no explicit $t$-dependence (scleronomic transformation) so $T=T_2$ is homogeneous quadratic in $\dot{\vec q}$
    \item Rheonomic transformations (a prescribed $\theta(t)=\omega_0t$, a moving support, \dots) generally give $T_1,T_0\neq0$, so $H\neq T+V$ even when $H$ is conserved
    \item $H$ is conserved whenever $\partial L/\partial t=0$ (no explicit $t$ left \emph{after} substituting any prescribed motion), regardless of whether $H=T+V$
\end{itemize}

\subsection*{Hamilton's canonical equations}
\begin{keybox}
\[
    \dot q_j = \frac{\partial H}{\partial p_j}, \qquad \dot p_j = -\frac{\partial H}{\partial q_j}
\]
\end{keybox}
$2M$ first-order ODEs replace the $M$ second-order Euler--Lagrange equations --- the natural form for numerical IVP integration. If $q_k$ is cyclic ($\partial H/\partial q_k=0$): $\dot p_k=0\Rightarrow p_k=$ const (matches the Lagrangian/Routhian result, $p_k\equiv B_k$).

\subsection*{Canonical transformations \& generating functions}
$(\vec q,\vec p)\to(\vec Q,\vec P)$ is \textbf{canonical} if Hamilton's equations hold in the new variables for some $H'(\vec Q,\vec P,t)$. A \textbf{type-1 generating function} $F(\vec q,\vec Q,t)$ produces one via
\begin{keybox}
\[
    p_j = \frac{\partial F}{\partial q_j}, \qquad P_j = -\frac{\partial F}{\partial Q_j}, \qquad H' = H+\frac{\partial F}{\partial t}
\]
\end{keybox}
\begin{itemize}[nosep, leftmargin=1.1em]
    \item If $F$ has no explicit $t$: $H'(\vec Q,\vec P)$ is just $H$ re-expressed in the new variables
    \item \textbf{Goal:} pick $F$ so $H'$ is simple (e.g.\ every $Q_j$ cyclic) --- then $\dot P_j=0$ and $\dot Q_j=\partial H'/\partial P_j=$ const, trivially integrable
    \item \textbf{Recipe:} (1) $p(q,Q)=\partial F/\partial q$; (2) $P(q,Q)=-\partial F/\partial Q$; (3) invert for $q(Q,P)$, then $p(Q,P)$; (4) substitute into $H(q,p)$ to get $H'(Q,P)$; (5) solve the (now trivial) EOM in $(Q,P)$; (6) transform back for $q(t),p(t)$
\end{itemize}

% ====================================================
% 1-D DYNAMICAL SYSTEMS & BIFURCATIONS
% ====================================================
\section*{1-D Dynamical Systems \& Bifurcations}
\subsection*{Fixed points \& linear stability}
System $\dot x=f(x;r)$, $r$ a control (parameter).
\begin{keybox}
\[
    \text{Fixed pts: } f(x^*;r)=0
    \qquad
    \text{Stability: } \dot\eta = f'(x^*)\,\eta,\ \ \eta\equiv x-x^*
\]
\end{keybox}
\begin{itemize}[nosep, leftmargin=1.1em]
    \item $f'(x^*)<0$: stable/attracting \quad $f'(x^*)>0$: unstable/repelling
    \item $f'(x^*)=0$: linearization inconclusive --- the \emph{marginal/degenerate} case where bifurcations live
\end{itemize}

\subsection*{Recipe: locate a bifurcation}
\begin{enumerate}[nosep, leftmargin=1.3em]
    \item Find all fixed points $x^*(r)$: solve $f(x;r)=0$
    \item Compute $f_x(x;r)=\partial f/\partial x$ and evaluate along each branch $x^*(r)$
    \item A bifurcation sits at $(x_c,r_c)$ where the branch is \emph{degenerate}:
    \[
        f(x_c;r_c)=0 \quad\text{and}\quad f_x(x_c;r_c)=0
    \]
    (solve these two equations simultaneously for $x_c,r_c$ --- this is exactly where two branches of $x^*(r)$ meet, or a branch's stability flips)
    \item Classify using the recipe/table below (Taylor-expand $f$ about $(x_c,r_c)$ if it's not already in normal form)
\end{enumerate}

\subsection*{Fast decision path (given explicit $f(x,r)$)}
\begin{enumerate}[nosep, leftmargin=1.3em]
    \item Is $x=x_0$ a root of $f$ for \emph{every} $r$ (a ``trivial'' branch, often $x_0=0$)?\\
    \textbf{No} $\to$ solve $f=0,\ f_x=0$ together for isolated $(x_c,r_c)$ $\to$ \textbf{saddle-node}.\\
    \textbf{Yes} $\to$ continue.
    \item Is $f(x,r)$ odd about $x_0$, i.e.\ $f(x_0-x,r)=-f(x_0+x,r)$?\\
    \textbf{No} $\to$ \textbf{transcritical} at $r_c$ solving $f_x(x_0,r_c)=0$.\\
    \textbf{Yes} $\to$ \textbf{pitchfork} at that $r_c$; Taylor/series-expand $f\approx (r-r_c)(x-x_0) + a\,(x-x_0)^3$ (expand any $\sinh$, $\tfrac{x}{1+x^2}$, etc.\ to cubic order) --- $a<0$ supercritical, $a>0$ subcritical.
\end{enumerate}

\subsection*{Classification (normal forms)}
Let $\xi=x-x_c$, $\mu=r-r_c$.
\begin{itemize}[nosep, leftmargin=1.1em]
    \item \textbf{Saddle-node (fold):} $\dot\xi=\mu-\xi^2$ (or $\mu+\xi^2$) --- \emph{no} fixed point at $\xi=0$ for all $\mu$; two fixed points ($x^*\propto\pm\sqrt\mu$) collide and annihilate as $\mu$ crosses $0$, existing only on one side of $r_c$. Test: $f_r\neq0,\ f_{xx}\neq0$.
    \item \textbf{Transcritical:} $\dot\xi=\mu\xi-\xi^2$ --- $\xi=0$ \emph{is} a fixed point for \emph{all} $\mu$ (trivial branch persists forever); a second branch $\xi=\mu$ crosses it at $\mu=0$ and the two \emph{exchange stability}; total \# of fixed points is unchanged across $r_c$. Test: $f_{x\mu}\neq0,\ f_{xx}\neq0$.
    \item \textbf{Supercritical pitchfork:} $\dot\xi=\mu\xi-\xi^3$ --- symmetric ($f(-\xi)=-f(\xi)$); trivial branch $\xi=0$ stable for $\mu<0$, goes unstable at $\mu=0$ and sheds \emph{two new stable} branches $\xi=\pm\sqrt\mu$ for $\mu>0$. Cubic coefficient $<0$.
    \item \textbf{Subcritical pitchfork:} $\dot\xi=\mu\xi+\xi^3$ --- same symmetry, but sheds \emph{two unstable} branches for $\mu<0$; trivial branch is unstable for $\mu>0$ with no nearby stable state (real systems jump/hysterese to a distant attractor via a higher-order saddle-node). Cubic coefficient $>0$.
\end{itemize}

\subsection*{Phase line (1-D phase portrait)}
Plot $\dot x$ vs.\ $x$ for fixed $r$: mark equilibria where the curve crosses $\dot x=0$, arrows on the $x$-axis between crossings showing the sign of $\dot x$ (right if $f>0$, left if $f<0$).
\begin{itemize}[nosep, leftmargin=1.1em]
    \item Filled dot: stable ($f$ goes $+\to-$ through it, arrows point \emph{in}); open dot ($\circ$): unstable ($-\to+$, arrows point \emph{out})
    \item Half-filled dot: semi-stable (same sign on both sides --- the signature of a saddle-node exactly at $r=r_c$)
    \item Sweeping $r$ and stacking phase lines side-by-side (equilibria vs.\ $r$) reconstructs the bifurcation diagram; solid branch = stable, dashed = unstable
\end{itemize}
\emph{(2-D extension --- nullclines, Jacobian trace/det, Hopf bifurcation: TBD)}

% ====================================================
% NONDIMENSIONALIZATION
% ====================================================
\section*{Nondimensionalizing an Equation of Motion}
\subsection*{Recipe}
\begin{enumerate}[nosep, leftmargin=1.3em]
    \item Write the dimensional EOM and list every dimensional quantity appearing in it (state variable(s), $t$, and all parameters: $m,k,b,g,\ell,\dots$)
    \item Choose a characteristic \textbf{length} $L$ and \textbf{time} $T$ built from the problem's own parameters (e.g.\ $L=$ a natural length/amplitude in the geometry, $T=\sqrt{m/k}$ or $1/\omega_n$) --- these are what get scaled out
    \item Define dimensionless variables $\xi=x/L$, $\tau=t/T$ (one pair per independent physical dimension in play; add a scaled force/energy if one appears)
    \item Convert derivatives by chain rule: $\dot x=\dfrac{L}{T}\xi'$, $\ddot x=\dfrac{L}{T^{2}}\xi''$ (primes $=d/d\tau$)
    \item Substitute into the EOM; divide the whole equation by the coefficient of the \emph{highest} derivative so it has unit leading coefficient
    \item Collect every remaining combination of parameters into as few dimensionless groups as possible (Buckingham $\pi$: \#groups $=$ \#parameters $-$ \#independent dimensions); name each one ($\zeta,\epsilon,\dots$) --- these are the only ``knobs'' left in the problem
    \item Sanity check: every leftover coefficient multiplying a $\xi,\xi',\xi''$ term must be a pure number (no residual kg/m/s)
\end{enumerate}

\subsection*{Worked template (spring--mass--damper type)}
Dimensional: $m\ddot x+b\dot x+kx=F(t)$ (e.g.\ a gravity component or drive). With $T=\sqrt{m/k}$, $L=$ a natural length scale, $\tau=t/T$, $\xi=x/L$:
\begin{keybox}
\[
    \xi''+2\zeta\,\xi'+\xi=\tilde F(\tau), \qquad
    2\zeta\equiv\frac{b}{\sqrt{mk}}, \qquad
    \tilde F\equiv\frac{F\,T^{2}}{mL}
\]
\end{keybox}
$\zeta$ (damping ratio) is the one leftover dimensionless group here; a tilted/geometric system typically adds one more (an angle $\alpha$ --- already dimensionless --- or a length ratio $a/L$), matching ``3 dimensionless parameters'' style results.

\subsection*{Overdamped/singular limit $\Rightarrow$ drops the order by one}
If a derivative's dimensionless coefficient is formally small (heavy damping $\Rightarrow$ small effective inertia, or a genuinely small mass), write it as $\epsilon\,\xi''+\xi'+\dots=0$ with $\epsilon\ll1$ and set $\epsilon\to0$ (quasi-static/regular-perturbation assumption):
\begin{keybox}
\[
    \epsilon\,\xi''+\xi' = g(\xi;\vec p) \ \ \xrightarrow{\ \epsilon\to0\ }\ \ \xi' = g(\xi;\vec p)
\]
\end{keybox}
This collapses a 2nd-order (2-D phase-plane) system to a \textbf{1-D, 1st-order} system --- the phase-line/bifurcation analysis above applies directly, with the remaining dimensionless groups $\vec p$ playing the role of $r$.

\end{multicols}
\end{document}
